\documentclass[11pt]{article}
\usepackage[T1]{fontenc}
\usepackage[utf8]{inputenc}
\usepackage{lmodern}
\usepackage[margin=1in]{geometry}
\usepackage{amsmath,amssymb,amsthm,mathtools}
\usepackage{microtype}
\usepackage[colorlinks=true,linkcolor=blue,citecolor=blue,urlcolor=blue]{hyperref}
\newtheorem{theorem}{Theorem}[section]
\newtheorem{proposition}[theorem]{Proposition}
\newtheorem{lemma}[theorem]{Lemma}
\newtheorem{corollary}[theorem]{Corollary}
\theoremstyle{definition}
\newtheorem{definition}[theorem]{Definition}
\theoremstyle{remark}
\newtheorem{remark}[theorem]{Remark}
\newcommand{\Fp}{\mathbb F_p}
\newcommand{\Fq}{\mathbb F_q}
\newcommand{\Tr}{\operatorname{Tr}}
\newcommand{\Nrm}{\operatorname{N}}
\newcommand{\id}{\operatorname{id}}
\newcommand{\rank}{\operatorname{rank}}
\DeclareMathOperator{\lcm}{lcm}
\numberwithin{equation}{section}
\title{Classification of Costas polynomials over finite extension fields}
\author{Kaiwen Jin\\ \small\href{mailto:yui.ui.siki@gmail.com}{yui.ui.siki@gmail.com}}
\date{November 2026}
\begin{document}
\maketitle
\begin{abstract}
A Costas polynomial over a finite field fixes zero and has a permutation
as every multiplicative difference \(f(dx)-f(x)\), \(d\ne1\).
We prove that over every nontrivial finite extension of a prime field,
these functions are exactly the permutation monomials followed by
invertible prime-field-linear maps. This establishes the extension-field
classification conjectured by Muratovi\'c-Ribi\'c, Pott, Thomson, and Wang,
without a shifting hypothesis. The proof combines the difference-set
structure of scalar trace components with a count of their Frobenius
multiplier parameters. Two rank inequalities force a full-length
Frobenius orbit, allowing a normalization that leaves only norm-dependent
terms. A direction theorem and a three-root multiplicative character
estimate then give the classification.
\end{abstract}

\section{Introduction}
\label{sec:introduction}

Costas polynomials describe a distinct-difference property in which
the argument is shifted multiplicatively and the values are compared
additively. More precisely, over a finite field \(\Fq\), a function
\(f:\Fq\to\Fq\) with \(f(0)=0\) is Costas if
\[
 x\longmapsto f(dx)-f(x)
\]
is a permutation for every \(d\in\Fq\setminus\{1\}\).
Every function on \(\Fq\) has a unique polynomial representative of
degree less than \(q\), so the terminology imposes no additional
polynomial restriction. These functions arise from circular Costas
structures and are related to relative difference sets and finite
geometries~\cite{Costas2015,RubioTorres}. Their permutation differences
also produce complete families of mutually orthogonal Latin
squares~\cite{Costas2026}.

The basic construction is
\[
 f(x)=L(x^s),\qquad \gcd(s,q-1)=1,
\]
where \(L\) is an invertible linear map over the prime field.
Muratovi\'c-Ribi\'c, Pott, Thomson, and Wang proved the prime-field
classification and conjectured that this construction gives all
Costas polynomials over extension fields~\cite[author version,
Corollary~4.2 and Conjecture~4.4]{Costas2015}.
Torres and Rubio proved the extension-field classification under the
stronger shifting condition: for each \(d\ne1\), there is an \(a\ne0\)
such that \(f(dx)-f(x)=f(ax)\) for all \(x\)
\cite[Definition~35 and Theorem~38 of arXiv:2210.16661v1]{RubioTorres}.
The orthomorphism formulation of
Muratovi\'c-Ribi\'c and Bala\v sev-Samarski retains the unrestricted
classification as a conjecture~\cite[Sections~1--2]{Costas2026}.
The issue is whether bijectivity of the differences alone forces the
same form.

\begin{theorem}\label{thm:main}
Let \(p\) be a prime, \(r\ge2\), \(q=p^r\), and let
\(f:\Fq\to\Fq\) satisfy \(f(0)=0\).
Then \(x\mapsto f(dx)-f(x)\) is a permutation of \(\Fq\) for every
\(d\in\Fq\setminus\{1\}\) if and only if there exist a positive
integer \(s\) with \(\gcd(s,q-1)=1\) and an invertible
\(\Fp\)-linear map \(L:\Fq\to\Fq\) such that
\[
 f(x)=L(x^s)\qquad(x\in\Fq).
\]
Equivalently, \(L\) is a linearized permutation polynomial
\(L(X)=\sum_{j=0}^{r-1}a_jX^{p^j}\).
\end{theorem}

The theorem has the same function class and conclusion as the
extension-field conjecture.

The main difficulty is to recover one vector-valued structure from all
its scalar components. Each nonzero trace component is balanced and
difference-balanced. Results of Pott and Wang give homogeneity and
relative difference sets for these scalar functions, and multiplier
theorems give a Frobenius multiplier for each component.
Counting the fixed points of \(x\mapsto bx^p\) shows that these
multiplier parameters run bijectively through a multiplicative subgroup
as the trace functional runs through the projective dual space.
Comparing this bijection with the image and kernel of a trace
coefficient gives two complementary rank inequalities.

Those inequalities force each active exponent orbit to have a rank
determined by its common divisor with \(q-1\).
If all ranks were smaller than \(r\), a cyclotomic divisor would give
a nontrivial multiplicative period of \(f\). Thus one full-rank orbit
can be changed to the term \(x\). In this normalization, coefficient
comparison leaves only terms controlled by the norm to the prime
field. Characteristic two then follows at once. In odd characteristic
the remaining possibilities are a binomial, excluded by Ball's
direction theorem, or scalar multiplication on multiplicative cosets.
For the latter, a three-root character estimate forces a collision
unless the successive scalars have a constant ratio. That condition
gives a permutation monomial.

Section~\ref{sec:multipliers} establishes the scalar multiplier
properties. Section~\ref{sec:coefficients} proves the rank identities
and the normalized form. Section~\ref{sec:completion} completes the
two remaining cases and proves Theorem~\ref{thm:main}.

\section{Scalar components and Frobenius multipliers}
\label{sec:multipliers}

Throughout, let \(p\) be a prime, \(r\geq2\), and
\[
 K=\Fq,\qquad q=p^r,\qquad m=q-1,\qquad
 n=\frac{m}{p-1}.
\]
Write \(K^*=K\setminus\{0\}\), and put
\[
 T(z)=\Tr_{K/\Fp}(z)=\sum_{j=0}^{r-1}z^{p^j},
 \qquad U=(K^*)^{p-1}.
\]
The group \(K^*\) is cyclic of order \(m\), and \(U\) has order \(n\).
All polynomial representatives on \(K\) have degree less than \(q\).
On \(K^*\), exponents are read modulo \(m\).

Call \(u:K\to K\) a Costas function if \(u(0)=0\) and
\begin{equation}\label{eq:costas}
 x\longmapsto u(dx)-u(x)
 \quad\hbox{is bijective for every }d\in K\setminus\{1\}.
\end{equation}
In particular, \(d=0\) shows that \(u\) is a permutation. For
\(\lambda\in K\), define its scalar component by
\[
 h_\lambda(x)=T(\lambda u(x)).
\]
We use the nondegeneracy of the trace pairing: every \(\Fp\)-linear
functional on \(K\) is \(z\mapsto T(\lambda z)\) for a unique \(\lambda\).
These functionals separate points.

\subsection{Balance and prime-field homogeneity}

\begin{lemma}\label{lem:balance}
If \(u\) is Costas and \(\lambda\ne0\), then \(h_\lambda\) and every
difference \(h_\lambda(dx)-h_\lambda(x)\), \(d\ne1\), take each value in
\(\Fp\) exactly \(p^{r-1}\) times on \(K\).
On \(K^*\), each has \(p^{r-1}-1\) zeros and \(p^{r-1}\) occurrences of
each nonzero value. Moreover, \(h_\lambda\) has no multiplicative period
other than \(1\).
\end{lemma}
\begin{proof}
The nonzero trace functional has fibers of size \(p^{r-1}\).
Apply it to the permutations \(u\) and \(u(d\,\cdot)-u\).
Removing the input \(0\), where all these functions vanish, gives the
counts on \(K^*\). If \(h_\lambda(ax)=h_\lambda(x)\) for every \(x\) and
\(a\ne1\), the corresponding difference is identically zero, contradicting
the same counts.
\end{proof}

For a scalar function on \(K^*\), we use \emph{balanced} to mean exactly
the deficient-zero counts in Lemma~\ref{lem:balance}. It is
\emph{difference-balanced} if all its differences with multipliers
in \(K^*\setminus\{1\}\) have those counts. The normalization matters:
no constant is added to the scalar components.

We use two results of Pott and Wang~\cite[Theorem~5.4 and
Corollary~5.5]{PottWang}, in the normalization of their Remark~3.2.
For odd \(p\), a balanced difference-balanced function
\(h:K^*\to\Fp\) satisfies
\begin{equation}\label{eq:scalarhom}
 h(ax)=a^d h(x)\qquad(a\in\Fp^*,\ x\in K^*)
\end{equation}
for some \(d\) coprime to \(p-1\).
Its nonzero fibers
\(R_b=\{x\in K^*:h(x)=b\}\) are relative difference sets in \(K^*\),
relative to \(\Fp^*\), with parameters
\begin{equation}\label{eq:rdsparams}
 (n,p-1,p^{r-1},p^{r-2}).
\end{equation}
Their images under \(\pi:K^*\to K^*/\Fp^*\) are ordinary difference
sets with parameters
\[
 (n,p^{r-1},(p-1)p^{r-2}),
\]
and their complement is \(\pi(R_0)\).

Here a \((v,k,\ell)\) difference set in a group of order \(v\) is a
\(k\)-element set in which every nonidentity group element has exactly
\(\ell\) ordered representations as a quotient of two set elements.
A relative \((v,w,k,\ell)\) difference set has ambient group order \(vw\)
and a specified subgroup of order \(w\); quotients outside that subgroup
have \(\ell\) representations, and its nonidentity elements have none.
The projection is then an ordinary \((v,k,w\ell)\) difference set.

\begin{lemma}\label{lem:homogeneity}
For every Costas function \(u\), there is a positive integer \(d\) with
\(\gcd(d,p-1)=1\) such that
\begin{equation}\label{eq:commonhom}
 u(ax)=a^d u(x)\qquad(a\in\Fp^*,\ x\in K).
\end{equation}
The polynomial representatives of \(u\) and all \(h_\lambda\) have
zero constant coefficient and zero coefficient at \(X^m\).
Every exponent with a nonzero coefficient is congruent to \(d\) modulo
\(p-1\), and hence is coprime to \(p-1\).
\end{lemma}
\begin{proof}
Suppose first that \(p\) is odd. By Lemma~\ref{lem:balance} and
\eqref{eq:scalarhom}, each nonzero component is homogeneous with an
exponent coprime to \(p-1\).
Choose components \(h_1,\ldots,h_r\) corresponding to a basis of the
dual space of \(K\). They are linearly independent as functions, since
\(u\) is surjective. The component \(h_i+h_j\), for \(i\ne j\), is
also nonzero. If the three homogeneity exponents are \(d_i,d_j,d_{ij}\),
then, for every \(a\in\Fp^*\),
\[
 (a^{d_i}-a^{d_{ij}})h_i+(a^{d_j}-a^{d_{ij}})h_j=0.
\]
Independence and cyclicity of \(\Fp^*\) give
\(d_i\equiv d_j\equiv d_{ij}\pmod{p-1}\).
Thus the exponents have a common residue \(d\). The dual basis
separates points, giving \eqref{eq:commonhom}.
For \(p=2\) the same identity holds with \(d=1\).

All constant coefficients vanish because \(u(0)=h_\lambda(0)=0\).
The scalar balance counts, together with \(r\ge2\), give
\(\sum_{x\in K}h_\lambda(x)=0\) in \(\Fp\); trace separation gives
\(\sum_{x\in K}u(x)=0\).
For \(1\le j<m\), the sum of \(x^j\) over \(K\) is zero, while
\(\sum_{x\in K}x^m=-1\). Hence the coefficients at \(X^m\) vanish.
Finally compare coefficients in \eqref{eq:commonhom}, and in its scalar
versions. Every active exponent is \(d\) modulo \(p-1\).
\end{proof}

\subsection{A common multiplier for the scalar fibers}

A numerical multiplier \(t\) of a subset \(D\) of a finite abelian
group means that powering by \(t\) is a group automorphism and
\(D^{(t)}=aD\) for some group element \(a\).
We use Hall's sufficient multiplier theorem in the form stated by
Turyn~\cite[p.~386]{Turyn}: for a cyclic \((v,k,\ell)\) difference set,
let \(M\) be a positive divisor of \(k-\ell\), with
\[
 M>\ell,\qquad \gcd(M,v)=1.
\]
An integer \(t\) coprime to \(v\) is a multiplier if, for every prime
\(P\mid M\), some nonnegative power of \(P\) is congruent to \(t\)
modulo \(v\).

We also use the following specialization of
Pott--Reuschling--Schmidt~\cite[Result~2]{PRS}.
Let \(R\) be an abelian relative \((v,w,p^a,\ell)\) difference set
in a group of exponent \(e\), with \(a\ge1\), \(p\nmid e\), and \(p^a>\ell\).
If \(p\) is a multiplier of the projected ordinary difference set,
then \(p\) is a multiplier of \(R\).
In Result~2, choose the divisor of the relative set's size to be
\(k_1=p^a\). Its second size parameter is
\(k_2=k_1/\gcd(e,k_1)=p^a\), so \(k_2>\ell\).
The only prime factor \(P=p\) is coprime to \(e\), allowing its
auxiliary integer to be \(Q=p\).
Taking the identity as an already existing multiplier and the
auxiliary exponent to be \(1\) gives the required congruence
\(1\cdot Q^1\equiv p\pmod e\).

\begin{proposition}\label{prop:scalar-multiplier}
For every Costas function \(u\) and every \(\lambda\ne0\), there is a
unique \(\gamma_\lambda\in K^*\) such that
\begin{equation}\label{eq:scalar-multiplier}
 h_\lambda(x^p)=h_\lambda(\gamma_\lambda x)\qquad(x\in K).
\end{equation}
Both \(\gamma_\lambda\) and
\(\beta_\lambda=\gamma_\lambda^{-p}\) are constant on each
nonzero \(\Fp\)-line in the parameter \(\lambda\).
\end{proposition}
\begin{proof}
We first show that one \(a\in K^*\) satisfies
\begin{equation}\label{eq:fiber-multiplier}
 R_b^{(p)}=aR_b\qquad(b\in\Fp),
 \qquad R_b=\{x\in K^*:h_\lambda(x)=b\}.
\end{equation}

If \(p=2\), put \(Z=R_0\). For a multiplier \(c\ne1\), the two
off-diagonal counts of \((h_\lambda(x),h_\lambda(cx))\) are equal,
since both coordinates have the same number of zeros. Their sum is
\(q/2\) by difference balance. The zero-zero count is therefore
\(q/4-1\). Thus \(Z\) is a cyclic
\((q-1,q/2-1,q/4-1)\) difference set.
When \(r\ge3\), Hall's theorem applies with \(M=q/4\) and \(t=2\):
\(M\) is a power of \(2\), is coprime to \(q-1\), and exceeds
\(q/4-1\). When \(r=2\), \(Z\) is a singleton, and the multiplier
assertion is direct. Since \(R_1\) is its complement in \(K^*\),
the same translate works for both fibers.

Now suppose \(p\) is odd. By \eqref{eq:rdsparams},
\(C=\pi(R_0)\) is the complement of the projected nonzero fiber.
The complement parameters are
\[
 \left(n,\frac{p^{r-1}-1}{p-1},
            \frac{p^{r-2}-1}{p-1}\right).
\]
For \(r\ge3\), its difference-set order is \(p^{r-2}\), which is
coprime to \(n\) and exceeds its last parameter.
Hall's theorem with \(M=p^{r-2}\), \(t=p\), shows that \(p\) is a
multiplier of \(C\). For \(r=2\), \(C\) is a singleton.
Complementation therefore shows in both cases that \(p\) is a
multiplier of \(\pi(R_1)\).
The relative multiplier result applies to \(R_1\), since its size
is \(p^{r-1}>p^{r-2}\) and \(p\nmid m\).
Consequently \(R_1^{(p)}=aR_1\) for some \(a\).
By Lemma~\ref{lem:homogeneity}, every nonzero fiber is \(cR_1\)
for a suitable \(c\in\Fp^*\). Frobenius fixes \(c\), so the same
\(a\) works for all nonzero fibers. It works for \(R_0\) as well,
by taking complements. This proves \eqref{eq:fiber-multiplier}.

Let \(\gamma_\lambda^p=a\). Equation~\eqref{eq:fiber-multiplier}
is equivalent to
\(h_\lambda(a^{-1}y^p)=h_\lambda(y)\) on \(K^*\).
Substitute \(y=\gamma_\lambda x\) to obtain
\eqref{eq:scalar-multiplier}; it also holds at \(0\).
Two choices of \(\gamma_\lambda\) would give a nontrivial multiplicative
period of \(h_\lambda\), excluded by Lemma~\ref{lem:balance}.
For \(c\in\Fp^*\), \(h_{c\lambda}=c h_\lambda\), so uniqueness
also gives the asserted constancy on parameter lines.
\end{proof}

\subsection{Counting the multiplier parameters}

\begin{proposition}\label{prop:projective}
The map \(\lambda\mapsto\beta_\lambda\) induces a bijection from
the nonzero \(\Fp\)-lines in \(K\) onto \(U\).
More precisely, for every \(b\in K^*\),
\begin{equation}\label{eq:beta-count}
 \#\{\lambda\in K^*:\beta_\lambda=b\}
 =\begin{cases}p-1,&b\in U,\\0,&b\notin U.\end{cases}
\end{equation}
\end{proposition}
\begin{proof}
Since \(u\) is injective, \(u(bx^p)=u(x)\) if and only if
\(bx^p=x\). The number \(C_b\) of such \(x\) is \(p\) when \(b\in U\)
and \(1\) otherwise: besides \(0\), the equation is
\(x^{p-1}=b^{-1}\).

Fix a primitive complex \(p\)-th root \(\zeta\).
Additive character orthogonality gives
\begin{equation}\label{eq:fixedorthog}
 qC_b=\sum_{\lambda\in K}\sum_{x\in K}
 \zeta^{\,h_\lambda(bx^p)-h_\lambda(x)}.
\end{equation}
The \(\lambda=0\) term is \(q\).
For \(\lambda\ne0\), choose \(v\) with \(v^p=b\).
By \eqref{eq:scalar-multiplier} the inner sum is
\[
 \sum_{x\in K}\zeta^{\,h_\lambda(\gamma_\lambda vx)-h_\lambda(x)}.
\]
It equals \(q\) if \(\gamma_\lambda v=1\), and equals \(0\) otherwise
by Lemma~\ref{lem:balance}. The first condition is precisely
\(b=\beta_\lambda\). Dividing \eqref{eq:fixedorthog} by \(q\) proves
\eqref{eq:beta-count}. Each nonzero parameter line has \(p-1\) elements,
and Proposition~\ref{prop:scalar-multiplier} makes \(\beta\) constant
on it. The fiber count proves the bijection.
\end{proof}

\section{Trace coefficients and normalization}
\label{sec:coefficients}

\subsection{Frobenius orbits and their ranks}

Let \(S\) contain the least positive representative of each orbit of
multiplication by \(p\) on the residues \(1,\ldots,m-1\).
For \(t\in S\), let \(\ell_t\) be its orbit length and put
\(E_t=\mathbb F_{p^{\ell_t}}\subset K\).
The length \(\ell_t\) divides \(r\), and \(x^t\in E_t\) for \(x\in K\).

\begin{lemma}\label{lem:trace-expansion}
For a Costas function \(u\), there are unique \(\Fp\)-linear maps
\(A_t:K\to E_t\), \(t\in S\), such that
\begin{equation}\label{eq:trace-expansion}
 h_\lambda(x)=\sum_{t\in S}
 \Tr_{E_t/\Fp}\bigl(A_t(\lambda)x^t\bigr)
 \qquad(\lambda,x\in K).
\end{equation}
There are unique \(\Fp\)-linear maps \(L_t:E_t\to K\) with
\begin{equation}\label{eq:adjoint}
 T(\lambda L_t(z))=\Tr_{E_t/\Fp}(A_t(\lambda)z),
 \qquad u(x)=\sum_{t\in S}L_t(x^t).
\end{equation}
The term \(L_t(x^t)\) has polynomial exponents only in the \(t\)-orbit.
It is nonzero exactly when \(A_t\) is nonzero.
\end{lemma}
\begin{proof}
Since \(h_\lambda\) is \(\Fp\)-valued, \(h_\lambda^p=h_\lambda\).
Lemma~\ref{lem:homogeneity} removes the constant and \(X^m\) terms.
Comparing coefficients on \(K^*\) shows that the coefficients along
the orbit of \(t\) are
\[
 A_t(\lambda),\ A_t(\lambda)^p,\ \ldots,\
 A_t(\lambda)^{p^{\ell_t-1}}.
\]
Closing the orbit gives \(A_t(\lambda)^{p^{\ell_t}}=A_t(\lambda)\).
Coefficient extraction is \(\Fp\)-linear in \(\lambda\), proving
\eqref{eq:trace-expansion} and uniqueness.

The nondegenerate trace pairings on \(K\) and \(E_t\) give the unique
adjoint \(L_t\) in \eqref{eq:adjoint}. Summing that identity and using
trace separation gives the expression for \(u\).
Every \(\Fp\)-linear endomorphism of \(K\) is uniquely a linearized
polynomial \(\sum_{j=0}^{r-1}a_jX^{p^j}\): evaluation is injective by
the degree bound \(p^{r-1}<q\), and both vector spaces have dimension
\(r^2\) over \(\Fp\). Extend \(L_t\) linearly from \(E_t\) to \(K\);
substituting \(X^t\) in such a polynomial shows the orbit assertion.
Trace nondegeneracy also gives \(L_t=0\) if and only if \(A_t=0\).
No division by an orbit length is involved.
\end{proof}

Call the \(t\)-orbit active when \(A_t\ne0\).
Taking the coefficient at exponent \(t\) in
\eqref{eq:scalar-multiplier} gives
\[
 A_t(\lambda)^{p^{\ell_t-1}}
       =A_t(\lambda)\gamma_\lambda^t.
\]
Raise this identity to the \(p\)-th power. For every \(\lambda\ne0\),
\begin{equation}\label{eq:coefficient-multiplier}
 A_t(\lambda)=0
 \quad\hbox{or}\quad
 A_t(\lambda)^{p-1}=\beta_\lambda^t.
\end{equation}

\begin{proposition}\label{prop:rank}
For an active orbit \(t\), let \(a=\rank_{\Fp} A_t\) and
\(c=\gcd(t,m)\). Then
\begin{equation}\label{eq:rank-result}
 a=\ell_t,\qquad a\mid r,\qquad
 c=\frac{m}{p^a-1}.
\end{equation}
In particular, \(A_t\) is onto \(E_t\). If \(t\) is coprime to \(m\),
then \(a=r\) and \(A_t:K\to K\) is invertible.
\end{proposition}
\begin{proof}
We have \(1\le a\le\ell_t\le r\).
An active scalar coefficient implies an active polynomial exponent of
\(u\) in the same orbit. By Lemma~\ref{lem:homogeneity},
\(\gcd(t,p-1)=1\), so
\(c=\gcd(t,m)=\gcd(t,n)\).
The image of the \(t\)-th power map on \(U\) has size \(n/c\).

The nonzero \(\Fp\)-lines in \(\operatorname{im}A_t\) give
\((p^a-1)/(p-1)\) distinct \((p-1)\)-st powers.
By \eqref{eq:coefficient-multiplier} and
Proposition~\ref{prop:projective}, all lie in \(U^t\). Hence
\begin{equation}\label{eq:rank-ineq-one}
 \frac{p^a-1}{p-1}\le\frac nc.
\end{equation}
For a fixed nonzero image value \(z\), its preimage under \(A_t\)
has \(p^{r-a}\) elements. They lie in distinct parameter lines:
if \(\lambda_2=b\lambda_1\) and both map to \(z\ne0\), then
\(bz=z\), so \(b=1\). The corresponding \(\beta_\lambda\) are
distinct elements of \(U\) satisfying
\(\beta_\lambda^t=z^{p-1}\).
Such a power-map fiber has \(c\) elements. Therefore
\begin{equation}\label{eq:rank-ineq-two}
 p^{r-a}\le c.
\end{equation}
Together these give the integer constraint
\[
 \frac{p^a-1}{p-1}\le\frac nc
 \le\frac{p^r-1}{(p-1)p^{r-a}}<\frac{p^a}{p-1}.
\]
The left endpoint is an integer and the interval has length at most
\(1\). Thus \(n/c=(p^a-1)/(p-1)\), yielding the formula for \(c\).
The identity
\(\gcd(p^a-1,p^r-1)=p^{\gcd(a,r)}-1\) gives \(a\mid r\).
Finally, \(p^j t\equiv t\pmod m\) is equivalent to
\(p^j\equiv1\pmod{p^a-1}\), whose least positive solution is \(a\).
The only degenerate modulus, \(p^a-1=1\), would force \(c=m\)
and \(t\equiv0\pmod m\), which is excluded. Thus \(\ell_t=a\).
\end{proof}

\begin{lemma}\label{lem:cyclotomic-divisor}
For \(r\ge2\),
\begin{equation}\label{eq:cyclotomic-divisor}
 \gcd_{\substack{a\mid r\\a<r}}\frac{p^r-1}{p^a-1}
 =\frac{p^r-1}{\lcm_{a\mid r,\ a<r}(p^a-1)}
 =\Phi_r(p)>1,
\end{equation}
where \(\Phi_r\) is the \(r\)-th cyclotomic polynomial.
\end{lemma}
\begin{proof}
All denominators divide \(p^r-1\), giving the first equality by prime
valuations. Every proper divisor of \(r\) divides \(r/P\) for some
prime \(P\mid r\); it suffices to take the least common multiple over
these maximal proper divisors.
For any finite list of prime valuations, the maximum is the
alternating sum of the minima over nonempty sublists. Using
\(\gcd(p^a-1,p^b-1)=p^{\gcd(a,b)}-1\), and writing
\(\mathcal P\) for the distinct prime divisors of \(r\), this gives
\[
 \frac{p^r-1}{\lcm_{a\mid r,\ a<r}(p^a-1)}
 =\prod_{J\subseteq\mathcal P}
       \bigl(p^{\,r/\prod_{P\in J}P}-1\bigr)^{(-1)^{|J|}}
 =\Phi_r(p).
\]
The last identity follows by grouping the roots of \(X^b-1\) by
their exact orders and applying inclusion-exclusion.
Each primitive \(r\)-th root of unity \(\xi\) differs from \(1\), so
\(|p-\xi|>p-1\ge1\). The positive integer \(\Phi_r(p)\), the product
of these factors, is therefore greater than \(1\).
\end{proof}

\begin{proposition}\label{prop:normalization}
Every Costas function has an active orbit \(t\) coprime to \(m\).
There are such a \(t\), its positive inverse \(e\) modulo \(m\), and
an invertible \(\Fp\)-linear map \(L_0:K\to K\), for which
\begin{equation}\label{eq:normalization}
 F(x)=L_0^{-1}\bigl(u(x^e)\bigr),\qquad
 u(x)=L_0\bigl(F(x^t)\bigr)
\end{equation}
is Costas and its entire exponent-\(1\) Frobenius orbit is exactly \(X\).
\end{proposition}
\begin{proof}
If every active orbit had rank a proper divisor \(a\) of \(r\),
Proposition~\ref{prop:rank} and Lemma~\ref{lem:cyclotomic-divisor}
would give a common divisor \(D=\Phi_r(p)>1\) of all active
exponents. Choose \(z\in K^*\) of order \(D\).
Then \eqref{eq:adjoint} gives \(u(zx)=u(x)\), contradicting
injectivity. Some active orbit therefore has rank \(r\), and
\eqref{eq:rank-result} makes its representative \(t\) coprime to \(m\).

For this orbit, the adjoint \(L_t\) in \eqref{eq:adjoint} is an
invertible map \(K\to K\); set \(L_0=L_t\).
If \(e>0\) satisfies \(et\equiv1\pmod m\), then
\eqref{eq:normalization} holds on all of \(K\).
Unit input powers and invertible \(\Fp\)-linear output maps preserve
the Costas property, since
\[
 J(u((dx)^e))-J(u(x^e))
       =J\bigl(u(d^e x^e)-u(x^e)\bigr)
\]
and \(d\ne1\) implies \(d^e\ne1\), also when \(d=0\).
The chosen term \(L_0(x^t)\) becomes \(x\).
Multiplication of exponents by \(e\) permutes their Frobenius orbits,
and a linearized output map moves exponents only within an orbit.
Thus no other orbit contributes to the exponent-\(1\) orbit of \(F\).
\end{proof}

We call a Costas function with this last polynomial property
\emph{normalized}. For such an \(F\), apply the preceding definitions
to \(F\). Its exponent-\(1\) coefficient satisfies
\(A_1(\lambda)=\lambda\). Equation~\eqref{eq:coefficient-multiplier}
then gives
\begin{equation}\label{eq:normalized-beta}
 \beta_\lambda=\lambda^{p-1}\qquad(\lambda\ne0).
\end{equation}

\subsection{The form of a normalized function}

\begin{lemma}\label{lem:power-differences}
The nonzero differences of distinct residues in
\(\{1,p,\ldots,p^{r-1}\}\) modulo \(p^r-1\) have unique ordered
representations.
\end{lemma}
\begin{proof}
Write \(a=b+L\pmod r\), with \(1\le L<r\).
The least nonnegative representative of \(p^a-p^b\) is the length-\(r\)
base-\(p\) word with digit \(p-1\) in the cyclic positions
\(b,b+1,\ldots,b+L-1\), and zero elsewhere. This follows from
\(p^L-1=(p-1)(1+p+\cdots+p^{L-1})\), reducing \(p^r\) to \(1\).
The word is neither all zero nor all \(p-1\); its unique beginning and
ending recover \(b\) and \(a\). Uniqueness of base-\(p\) digits proves
the assertion, including \(p=2,r=2\).
\end{proof}

\begin{lemma}\label{lem:block-form}
For a normalized \(F\), the coefficient map \(B=A_t\) of an active
orbit satisfies
\begin{equation}\label{eq:B}
 B(\lambda)^p=\lambda^{(p-1)t}B(\lambda)\qquad(\lambda\in K^*).
\end{equation}
If \((p-1)t\not\equiv0\pmod m\), then
\[
 B(\lambda)=c\lambda^{p^j},\qquad
 c\in\Fp^*,\quad 0\le j<r,\quad t\equiv p^j\pmod n.
\]
Its field-valued contribution to \(F\) is \(cx^v\), with
\(v\equiv1\pmod n\).
If \((p-1)t\equiv0\pmod m\), then \(n\mid t\), \(\ell_t=1\),
and its contribution is \(\gamma x^t\) for some \(\gamma\in K\).
\end{lemma}
\begin{proof}
Equations~\eqref{eq:coefficient-multiplier} and
\eqref{eq:normalized-beta} give \eqref{eq:B}, including when
\(B(\lambda)=0\).
Write the linearized polynomial \(B(X)=\sum_j c_jX^{p^j}\).
Let \(\delta=(p-1)t\not\equiv0\pmod m\).
Comparing supports in \eqref{eq:B} on \(K^*\), each nonzero
coefficient \(c_j\) requires
\(\delta+p^j\equiv p^i\pmod m\) for an \(i\ne j\).
Lemma~\ref{lem:power-differences} permits only one such ordered pair.
Thus \(B(X)=cX^{p^j}\).
Comparing its coefficient and exponent gives \(c^p=c\) and
\(\delta+p^j\equiv p^{j+1}\pmod m\).
Cancel \(p-1\) to obtain \(t\equiv p^j\pmod n\).
The map \(B\) has rank \(r\), so Proposition~\ref{prop:rank} gives
\(E_t=K\). Frobenius invariance of the trace yields
\[
 T(c\lambda^{p^j}x^t)
    =T\bigl(\lambda c x^{\,t p^{r-j}}\bigr).
\]
By the adjoint identity, the contribution is \(cx^v\), with
\(v\equiv t p^{r-j}\equiv1\pmod n\).

If \(\delta=0\), then \(n\mid t\), \(pt\equiv t\pmod m\), and
\(\ell_t=1\). The map \(B:K\to\Fp\) is a trace functional
\(B(\lambda)=T(\lambda\gamma)\); its adjoint contribution is
\(\gamma x^t\), as \(x^t\in\Fp\).
\end{proof}

\begin{proposition}\label{prop:normal-form}
Put \(\Nrm(x)=x^n\), the norm from \(K\) to \(\Fp\).
A normalized function has, on \(K^*\), the form
\begin{equation}\label{eq:norm-form}
 F(x)=x\sum_{k=0}^{p-2}c_k\Nrm(x)^k
            +\sum_{d=1}^{p-2}\gamma_d\Nrm(x)^d,
 \qquad c_k\in\Fp,\quad \gamma_d\in K,\quad c_0=1.
\end{equation}
It is homogeneous of degree \(1\) over \(\Fp\), and
\begin{equation}\label{eq:norm-restrictions}
 c_k\ne0\ \Longrightarrow\ rk\equiv0\pmod{p-1},\qquad
 \gamma_d\ne0\ \Longrightarrow\ rd\equiv1\pmod{p-1}.
\end{equation}
Consequently:
\begin{enumerate}
\item If \(p=2\), then \(F(x)=x\) on \(K\).
\item If \(p\) is odd and \(\gcd(r,p-1)=1\), then
\begin{equation}\label{eq:binomial-form}
 F(x)=x+\gamma x^{nd}
\end{equation}
on \(K\), where \(1\le d\le p-2\) and \(rd\equiv1\pmod{p-1}\).
\item If \(g=\gcd(r,p-1)>1\), then
\begin{equation}\label{eq:cyclo-form}
 F(x)=xV(x^{m/g})\qquad(x\in K^*),
\end{equation}
where \(V\in\Fp[X]\), \(\deg V<g\), and \(V\) is nonzero at every
\(g\)-th root of unity. In this case \(g\mid n\) and
\((p-1)\mid m/g\).
\end{enumerate}
\end{proposition}
\begin{proof}
Lemma~\ref{lem:block-form} leaves only exponents
\(1+kn\), \(0\le k\le p-2\), with \(\Fp\)-coefficients, and
\(nd\), \(1\le d\le p-2\), with \(K\)-coefficients.
They are disjoint because \(n>1\). The coefficient at \(X\) is \(1\).
This gives \eqref{eq:norm-form}.
Lemma~\ref{lem:homogeneity}, applied to the active exponent \(1\),
makes the common homogeneity degree \(1\) modulo \(p-1\).
Since \(n\equiv r\pmod{p-1}\), coefficient comparison gives
\eqref{eq:norm-restrictions}.

For \(p=2\) the first sum contains only \(c_0=1\), and the second
is empty. If \(p\) is odd and \(\gcd(r,p-1)=1\), only \(k=0\)
and the single indicated residue \(d\) can remain.
For \(g>1\), the second congruence in
\eqref{eq:norm-restrictions} has no solution; the first requires
\(k=j(p-1)/g\), \(0\le j<g\). This gives \eqref{eq:cyclo-form}.
Here \(g\mid r\), \(p\equiv1\pmod g\), and thus
\(n=1+p+\cdots+p^{r-1}\equiv r\equiv0\pmod g\).
Every \(g\)-th root \(y\) is \(x^{m/g}\) for some \(x\in K^*\)
by cyclicity. If \(V(y)=0\), then \(F(x)=F(0)=0\), contrary to
injectivity. All the assertions include \(0\) where stated, since
\(F(0)=0\) and all displayed monomial exponents are positive.
\end{proof}

\section{Completion of the classification}
\label{sec:completion}

\subsection{The binomial case}

We first exclude the second term in \eqref{eq:binomial-form}.
The connection with graph directions is useful because the Costas
condition controls products of slopes, rather than individual slopes.

\begin{lemma}\label{lem:binomial}
Suppose \(p\) is odd, \(c\in K^*\), and \(1<k\le q-1\).
If \(x+cx^k\) is Costas, then \(k=p^j\) for some \(0\le j<r\).
\end{lemma}
\begin{proof}
Let
\[
 D=\left\{\frac{y^k-x^k}{y-x}:x,y\in K,\ x\ne y\right\},
 \qquad D^+=D\setminus\{0\},
\]
and define subsets of \(K^*\) by
\[
 H=(K^*)^{k-1},\qquad
 S^+=\left\{\frac{a^k-1}{a-1}:
                   a\in K,\ a\ne1,\ a^k\ne1\right\}.
\]
The value \(a=0\) gives \(1\in S^+\).
Writing a secant's endpoints as \(x,ax\), with \(x\ne0\), gives
\[
 D^+=HS^+,\qquad S^+D^+=D^+D^+,
\]
since \(H\) is a subgroup. Secants through \(0\) are included by
the same value \(a=0\).

For \(a\ne1\), divide the Costas difference of \(x+cx^k\) by
\(a-1\). For every \(s=(a^k-1)/(a-1)\ne0\), the resulting map
\(x+csx^k\) is a permutation. A collision at distinct \(x,y\)
would be equivalent to
\[
 \frac{y^k-x^k}{y-x}=-\frac1{cs}.
\]
It follows that
\[
 -c^{-1}\notin S^+D^+=D^+D^+.
\]
Thus \(D^+\) and \(-c^{-1}(D^+)^{-1}\) are disjoint, whence
\begin{equation}\label{eq:directions-bound}
 2|D^+|\le q-1,\qquad |D|\le\frac{q+1}{2}.
\end{equation}

We apply Ball's direction theorem~\cite[Theorem~1.1]{Ball} to
the graph of \(v(x)=x^k\).
In its notation, let \(s\) be the largest power of \(p\) dividing
the size of every nonempty graph intersection with a line whose slope
is determined by the graph. The theorem gives: either \(s=1\) and
\(|D|\ge(q+3)/2\); or \(\mathbb F_s\subset K\) and
\[
 q/s+1\le |D|\le(q-1)/(s-1);
\]
or \(s=q\) and \(|D|=1\). Moreover, for \(s>2\) the graph is
\(\mathbb F_s\)-linear. The first possibility contradicts
\eqref{eq:directions-bound}. In odd characteristic any remaining
\(s>1\) is greater than \(2\). Therefore either the graph is a
linear subspace through \((0,0)\), or it is a line through \((0,0)\).
In both cases \(v\) is additive.

For completeness, additivity of \(x^k\) forces \(k\) to be a
\(p\)-power. The polynomial
\((X+Y)^k-X^k-Y^k\) has degree less than \(q\) in each variable
and vanishes on \(K^2\), so all its coefficients vanish.
Write \(k=\sum_j d_jp^j\) in base \(p\), and let \(j\) be its
least nonzero digit position. In
\[
 (1+X)^k=\prod_i(1+X^{p^i})^{d_i}
\]
over \(\Fp\), the coefficient at \(X^{p^j}\) is \(d_j\ne0\).
It would be a proper binomial coefficient unless \(k=p^j\).
The bound \(k<q\) gives \(j<r\).
\end{proof}

\begin{corollary}\label{cor:coprime-case}
If \(p\) is odd and \(\gcd(r,p-1)=1\), every normalized Costas
function is the identity.
\end{corollary}
\begin{proof}
In \eqref{eq:binomial-form}, the exponent \(nd\) satisfies
\(1<nd<m\) and is not coprime to \(m\), since \(n>1\) divides
both \(nd\) and \(m\). If \(\gamma\ne0\),
Lemma~\ref{lem:binomial} would make \(nd\) a \(p\)-power,
which is coprime to \(m\). Hence \(\gamma=0\).
\end{proof}

\subsection{A character count}

The remaining form \eqref{eq:cyclo-form} acts by a prime-field scalar
on each multiplicative coset of index \(g\). We will show that unequal
successive scalar ratios force a collision in a Costas difference.
The required simultaneous coset conditions follow from the next count.

\begin{lemma}\label{lem:size}
If \(g>1\) divides both \(r\) and \(p-1\), then
\[
 \sqrt q>2(g-1)^2.
\]
\end{lemma}
\begin{proof}
For \(g=2\), \(p\ge3,r\ge2\) give \(\sqrt q\ge3>2\).
For \(g=3\), the prime \(p\equiv1\pmod3\) is at least \(7\),
so \(\sqrt q\ge7^{3/2}>8\).
For \(g=4\), \(\sqrt q\ge5^2>18\).
For \(g\ge5\), use \(p\ge g+1\) and \(r\ge g\).
The sequence
\[
 R(v)=\frac{(v+1)^{v/2}}{2(v-1)^2}\qquad(v\ge5)
\]
has \(R(5)=36\sqrt6/32>1\), and
\[
 \frac{R(v+1)}{R(v)}
 >\sqrt{v+2}\,\frac{(v-1)^2}{v^2}
 >\frac{32}{25}>1.
\]
Thus \(\sqrt q\ge(g+1)^{g/2}>2(g-1)^2\).
\end{proof}

\begin{lemma}\label{lem:character-count}
Let \(g>1\) divide \(r\) and \(p-1\), let \(\chi\) be a complex
multiplicative character of \(K^*\) of exact order \(g\), and let
\(\omega\) be a primitive \(g\)-th root of unity.
Suppose \(A,B\in K^*\) are distinct and
\(\chi(A)=\chi(B)=\omega\). For any \(g\)-th root
\(\eta\ne1\), the number \(N\) of \(b\in K\setminus\{0,A,B\}\) with
\begin{equation}\label{eq:count-conditions}
 \chi(b)=1,\qquad
 \chi\left(\frac{b-A}{b-B}\right)=\eta
\end{equation}
satisfies
\begin{equation}\label{eq:count-lower}
 N\ge\frac{q+1-2(g-1)^2\sqrt q}{g^2}>0.
\end{equation}
\end{lemma}
\begin{proof}
Extend each nontrivial multiplicative character by zero at \(0\).
For \(0\le j,k<g\), put
\[
 S_{jk}=\sum_{b\notin\{0,A,B\}}
          \chi(b)^j\chi\left(\frac{b-A}{b-B}\right)^k.
\]
Character orthogonality on the group of \(g\)-th roots gives
\begin{equation}\label{eq:character-orthog}
 g^2N=\sum_{j,k=0}^{g-1}\eta^{-k}S_{jk}.
\end{equation}
Here \(S_{00}=q-3\). For \(j\ne0\), the full nontrivial character
sum is zero; removing \(A,B\) gives
\(S_{j0}=-2\omega^j\).
For \(k\ne0\), the substitution \(y=(b-A)/(b-B)\) maps
\(K\setminus\{B\}\) bijectively to \(K\setminus\{1\}\).
Removing \(b=0\) removes the value
\(\chi(A/B)^k=1\), and the \(b=A\) term is zero.
Thus \(S_{0k}=-2\). Consequently the two nonprincipal axes in
\eqref{eq:character-orthog} each contribute \(2\):
\[
 \sum_{j=1}^{g-1}S_{j0}=2,\qquad
 \sum_{k=1}^{g-1}\eta^{-k}S_{0k}=2.
\]

For \(1\le j,k<g\), the remaining sum is
\[
 S_{jk}=\sum_{b\in K}
       \chi\bigl(b^j(b-A)^k(b-B)^{g-k}\bigr).
\]
The polynomial is monic and has exactly the three distinct roots
\(0,A,B\). Each root multiplicity is positive and less than \(g\),
so the polynomial is not a \(g\)-th power.
We use the Weil bound in the form
of~\cite[Theorem~2.1]{AGW}: if \(\chi\) has order \(g>1\) and a
monic polynomial \(P\) is not a \(g\)-th power, then
\[
 \left|\sum_{b\in K}\chi(P(b))\right|
       \le (D-1)\sqrt q,
\]
where \(D\) is the number of distinct roots in a splitting field.
Its hypotheses here give \(|S_{jk}|\le2\sqrt q\).
The \((g-1)^2\) mixed terms have weights of modulus \(1\).
Since the left side of \eqref{eq:character-orthog} is real,
the triangle inequality proves the first inequality in
\eqref{eq:count-lower}. Lemma~\ref{lem:size} makes its numerator
strictly positive.
\end{proof}

\subsection{Rigidity on the multiplicative cosets}

\begin{proposition}\label{prop:cyclotomic-rigidity}
Let \(g>1\) divide both \(r\) and \(p-1\).
Suppose \(u\) is Costas and
\[
 u(x)=xV(x^{m/g})\qquad(x\in K^*),
\]
where \(V\in\Fp[X]\) is nonzero on the \(g\)-th roots of unity.
Then
\[
 u(x)=a_0x^{\,1+km/g}\qquad(x\in K)
\]
for some \(a_0\in\Fp^*\) and \(0\le k<g\), and the exponent
\(1+km/g\) is coprime to \(m\).
\end{proposition}
\begin{proof}
Choose a primitive \(\alpha\in K^*\), a primitive complex
\(g\)-th root \(\omega\), and the character
\(\chi(\alpha^j)=\omega^j\).
Since \(g\mid r,p-1\), we have \(g\mid n\) and
\((p-1)\mid m/g\). Thus \(\chi\) is trivial on \(\Fp^*\).
The element \(\rho=\alpha^{m/g}\) is a primitive \(g\)-th root
in \(\Fp\). On the coset \(\chi(x)=\omega^i\), write
\[
 a_i=V(\rho^i)\in\Fp^*,\qquad u(x)=a_i x,
 \qquad c_i=\frac{a_{i+1}}{a_i},
\]
with indices modulo \(g\).
Because \(\chi(a_i)=1\), \(u\) preserves each coset.
Conjugate multiplication by \(\alpha\) through \(u\):
\[
 G=u\circ(\alpha\,\id)\circ u^{-1}.
\]
For \(x\ne0\) with \(\chi(x)=\omega^i\),
\begin{equation}\label{eq:G}
 G(x)=\alpha c_i x,\qquad
 \prod_{i=0}^{g-1}c_i=1,\qquad
 G^g(x)=\alpha^g x.
\end{equation}
These identities also hold at \(0\) where appropriate.
Moreover, \(G^m=\id\).

For every integer \(k\not\equiv0\pmod m\), the map \(G^k-\id\)
is a permutation, because
\[
 (G^k-\id)\circ u
      =u(\alpha^k\,\cdot)-u
\]
is a Costas difference. Negative iterates use the inverse permutation.
If \(b\in\ker\chi\), write \(b=\alpha^{gt}\) for an integer \(t\).
Equation~\eqref{eq:G} gives \(G^{gt}=b\,\id\), a linear scalar map.
Therefore
\begin{equation}\label{eq:kernel-difference}
 G-b\,\id=(b\,\id)\circ(G^{1-gt}-\id)
\end{equation}
is a permutation. Indeed \(1-gt\not\equiv0\pmod m\), since
\(g>1\) divides \(m\).

Suppose \(c_i\ne c_j\). Then \(i\ne j\), and
\[
 A=\alpha c_i,\qquad B=\alpha c_j,\qquad
 \eta=\omega^{j-i}
\]
satisfy the hypotheses of Lemma~\ref{lem:character-count}:
\(\chi(A)=\chi(B)=\omega\), \(A\ne B\), and \(\eta\ne1\).
Choose the \(b\in\ker\chi\setminus\{0,A,B\}\) supplied by that lemma.
On input cosets \(i,j\), the map \(G-b\,\id\) multiplies by
\(A-b,B-b\), respectively. Their output cosets have characters
\[
 \omega^i\chi(A-b),\qquad \omega^j\chi(B-b).
\]
They are equal by \eqref{eq:count-conditions}.
Each nonzero scalar multiplication maps an entire input coset onto
the corresponding output coset. The two distinct input cosets
therefore have overlapping images, contradicting
\eqref{eq:kernel-difference}. Thus all \(c_i\) are equal to one \(c\).

Their product is \(1\), so \(c^g=1\) and \(c=\rho^k\) for some
\(0\le k<g\). The recurrence gives \(a_i=a_0\rho^{ki}\).
For \(x\) in coset \(i\), \(x^{m/g}=\rho^i\), whence
\[
 u(x)=a_i x=a_0x^{\,1+km/g}.
\]
Both sides vanish at \(0\). Finally \(u\) is a permutation, so
the power map with exponent \(1+km/g\) permutes \(K^*\);
cyclicity makes this equivalent to coprimality with \(m\).
\end{proof}

\begin{proof}[Proof of Theorem~\ref{thm:main}]
The converse follows directly: if \(f(x)=L(x^s)\), where
\(L\) is invertible and \(\Fp\)-linear and \(\gcd(s,m)=1\), then
\[
 f(dx)-f(x)=L\bigl((d^s-1)x^s\bigr)
\]
is a permutation for every \(d\ne1\), including \(d=0\).

For the forward implication, apply
Proposition~\ref{prop:normalization} to the given \(f\), obtaining
\(f(x)=L_0(F(x^t))\) with \(\gcd(t,m)=1\) and \(F\) normalized.
If \(p=2\), Proposition~\ref{prop:normal-form} gives \(F=\id\).
If \(p\) is odd and \(\gcd(r,p-1)=1\),
Corollary~\ref{cor:coprime-case} gives \(F=\id\).
In the remaining case set \(g=\gcd(r,p-1)>1\).
Proposition~\ref{prop:normal-form} gives \eqref{eq:cyclo-form},
and Proposition~\ref{prop:cyclotomic-rigidity} gives
\(F(x)=a_0x^v\), with \(a_0\in\Fp^*\) and \(\gcd(v,m)=1\).
In every case there are a positive unit exponent \(v\) and an
invertible \(\Fp\)-linear map \(J\) such that \(F(x)=J(x^v)\).
Thus
\[
 f(x)=(L_0\circ J)(x^{tv})\qquad(x\in K).
\]
The exponent \(s=tv\) is positive and coprime to \(m\), and
\(L=L_0\circ J\) is an invertible \(\Fp\)-linear map, as required.
\end{proof}

\paragraph{AI assistance.}
OpenAI Codex assisted in developing the proof, drafting the manuscript,
and obtaining automated reader and referee feedback. These automated
reviews are AI reviews.

\begin{thebibliography}{9}

\bibitem{AGW}
A.~Akbary, D.~Ghioca, and Qiang Wang,
\emph{On permutation polynomials of prescribed shape},
Finite Fields Appl. \textbf{15} (2009), 195--206.
\href{https://doi.org/10.1016/j.ffa.2008.12.001}{doi:10.1016/j.ffa.2008.12.001}.
Author version, Theorem~2.1:
\url{https://people.math.carleton.ca/~wang/papers/distributionPP.pdf}.

\bibitem{Ball}
S.~Ball,
\emph{The number of directions determined by a function over a finite field},
author version, 15 September 2003, Theorem~1.1.
\url{https://web.mat.upc.edu/simeon.michael.ball/hassedirection.pdf}.

\bibitem{Costas2015}
A.~Muratovi\'c-Ribi\'c, A.~Pott, D.~Thomson, and Qiang Wang,
\emph{On the characterization of a semi-multiplicative analogue of planar
functions over finite fields},
in \emph{Topics in Finite Fields}, Contemp. Math. \textbf{632},
Amer. Math. Soc., 2015, 317--325.
\href{https://doi.org/10.1090/conm/632/12635}{doi:10.1090/conm/632/12635}.
\href{https://www.academia.edu/79639197/On_the_characterization_of_a_semi_multiplicative_analogue_of_planar_functions_over_finite_fields}{Author version consulted}
(Corollary~4.2 and Conjecture~4.4).

\bibitem{Costas2026}
A.~Muratovi\'c-Ribi\'c and A.~Bala\v sev-Samarski,
\emph{Note on the Equivalence of Costas Polynomials and Orthomorphisms},
arXiv:2606.07097v1, 2026.
\url{https://arxiv.org/abs/2606.07097v1}.

\bibitem{PRS}
A.~Pott, D.~Reuschling, and B.~Schmidt,
\emph{A multiplier theorem for projections of affine difference sets},
J. Statist. Plann. Inference \textbf{62} (1997), 63--67, Result~2.
\url{https://personal.ntu.edu.sg/bernhard/Publications/pub/AffineDSMult.pdf}.

\bibitem{PottWang}
A.~Pott and Qi Wang,
\emph{Difference balanced functions and their generalized difference sets},
arXiv:1309.7842v1, 2013.
\url{https://arxiv.org/abs/1309.7842v1}.

\bibitem{RubioTorres}
J.~Torres and I.~Rubio,
\emph{Circular Costas maps: a multidimensional analog of circular Costas sequences},
Cryptogr. Commun. \textbf{15} (2023), 941--958.
\href{https://doi.org/10.1007/s12095-023-00654-2}{doi:10.1007/s12095-023-00654-2}.
Preprint v1 (I.~Rubio and J.~Torres), used for the numbered locators:
\url{https://arxiv.org/abs/2210.16661v1}.

\bibitem{Turyn}
R.~J.~Turyn,
\emph{The multiplier theorem for difference sets},
Canad. J. Math. \textbf{16} (1964), 386--388.
\href{https://doi.org/10.4153/CJM-1964-039-1}{doi:10.4153/CJM-1964-039-1}.

\end{thebibliography}
\end{document}
